% ============================================================
%  Journal of Integer Sequences -- submission draft
%  Title:  Exceptional Cases and Statistical Structure of
%          Sequence A247975: a Computational Study of
%          Sun's Conjecture 4.1(i)
%  Author: Carlo Corti
% ============================================================

\documentclass[12pt]{article}

\usepackage{amsmath,amssymb,amsthm}
\usepackage{hyperref}
\usepackage{booktabs}
\usepackage{geometry}
\geometry{margin=1in}

\newtheorem{conjecture}{Conjecture}
\newtheorem{theorem}{Theorem}
\newtheorem{proposition}{Proposition}
\newtheorem{remark}{Remark}

\title{Exceptional Cases and Statistical Structure of\\
	Sequence A247975: a Computational Study of\\
	Sun's Conjecture~4.1(i)}

\author{Carlo Corti\\[4pt]
	\small Independent researcher\\
	\small 21-040~Kalin\'owka, Poland\\
	\small \texttt{carlo.corti@outlook.com}}

\date{}

\begin{document}
	
	\maketitle
	
	\noindent
	\textit{Mathematics Subject Classification (2020):}
	Primary 11A41; Secondary 11Y55, 11N05.
	
	\noindent
	\textit{Keywords:}
	prime numbers, integer sequences, OEIS, computational number theory,
	Sun's conjecture, Hill estimator, Cram\'er ratio.
	
	% ----------------------------------------------------------
	%  Abstract
	% ----------------------------------------------------------
	\begin{abstract}
		We present a computational investigation of sequence A247975 in the
		OEIS, which encodes Sun's Conjecture~4.1(i): for every positive integer
		$m$, there exists a positive integer $n$ such that $m+n$ divides
		$p_m^2 + p_n^2$, where $p_k$ denotes the $k$-th prime.
		We determine $a(m)$ for $99\,951$ of the values $m = 1, 2, \ldots,
		100\,000$, with search bound $n \leq 2 \times 10^9$
		($p_n \leq 4.71 \times 10^{10}$); the remaining $49$ values are
		unresolved within this bound.
		Among the $99\,951$ resolved values we identify $398$ exceptional
		resolved cases where $a(m) > 10^8$, and document two values with
		outstanding Cram\'er-type ratio $a(m)/(m \log m)$
		(all logarithms are natural): $m = 4703$ with
		$a(m) = 760\,027\,770$ (ratio $\approx 19\,111$, the global maximum)
		and $m = 14\,740$
		with $a(m) = 1\,994\,463\,433$ (ratio $\approx 14\,097$, the largest
		absolute value of $a(m)$).
		Statistical analysis of the resolved dataset yields a Hill statistic
		$\hat{\alpha}_{\mathrm{Hill}} \approx 0.713$, consistent with
		heavy-tailed behaviour.
		The $49$ unresolved values all exhibit a filter-detected modular
		obstruction rate between $75.1\%$ and $77.3\%$.
		Divisibility was independently verified for $398$ exceptional cases via
		\texttt{SymPy}; minimality was independently verified for $290$ cases
		via an independent C sieve of Eratosthenes: $100$ exceptional cases
		($a(m) > 10^8$, the first such minimality check for this class),
		$150$ randomly selected non-exceptional cases, and all $40$ cases with
		$a(m) > 10^9$, including $m = 14\,740$ and $m = 16\,167$, the two largest
		resolved values.
		The computational work was developed with the assistance of the Claude
		AI language model (Anthropic); all computations were executed and
		verified by the author on a personal workstation.
	\end{abstract}
	
	% ----------------------------------------------------------
	\section{Introduction}
	% ----------------------------------------------------------
	
	In a paper submitted to arXiv in 2013 and published in 2019,
	Zhi-Wei Sun~\cite{Sun2019} proposed a collection of conjectures on
	addition and divisibility involving primes and other arithmetic
	functions.
	Among these, Conjecture~4.1(i) (p.~15) states the following.
	
	\begin{conjecture}[Sun, 2013/2019]\label{conj:sun}
		For every positive integer $m$, there exists a positive integer $n$
		such that
		\[
		(m + n) \mid p_m^2 + p_n^2,
		\]
		where $p_k$ denotes the $k$-th prime.
	\end{conjecture}
	
	The conjecture was submitted to the OEIS in September 2014 and appears
	as sequence A247975~\cite{OEIS_A247975}.
	The sequence $a(m)$, defined as the least positive integer $n$
	satisfying the divisibility condition, was computed for $m \leq 10\,000$
	by Chai Wah Wu~\cite{WuOEIS}.
	
	The present work has three main contributions.
	First, we extend the computation to $m = 1, \ldots, 100\,000$,
	providing the largest verified dataset for this conjecture to date.
	Second, we carry out a systematic statistical analysis of the
	distribution of $a(m)$.
	Third, we identify and characterize $49$ values of $m$ for which the
	conjecture resists verification within our search bound, all sharing a
	distinctive modular obstruction pattern.
	
	% ----------------------------------------------------------
	\section{Setup and Notation}
	% ----------------------------------------------------------
	
	We begin by reformulating the divisibility condition in modular
	arithmetic and establishing the number-theoretic constraints that
	underlie the computational search.
	
	For positive integers $m$ and $n$, set $s = m + n$.
	The divisibility condition $(m+n) \mid p_m^2 + p_n^2$ is equivalent to
	\[
	p_n^2 \equiv -p_m^2 \pmod{s}.
	\]
	When $\gcd(p_m, s) = 1$, this simplifies to
	\[
	\left( p_n p_m^{-1} \right)^2 \equiv -1 \pmod{s},
	\]
	so (under the assumption $\gcd(p_m, s) = 1$) a solution exists at
	index $n$ if and only if $p_n$ falls in one of
	the $2^{\omega(s)}$ residue classes of the form $r \cdot p_m \pmod{s}$,
	where $r$ ranges over the square roots of $-1$ modulo~$s$
	and $\omega(s)$ denotes the number of distinct prime factors of~$s$.
	
	When $\gcd(p_m, s) > 1$, i.e., $p_m \mid s = m+n$, the divisibility
	chain $p_m \mid s$, $p_m \mid p_m^2$, and $(m+n) \mid (p_m^2 + p_n^2)$
	together force $p_m \mid p_n^2$, hence $p_m \mid p_n$ (since $p_m$ is
	prime), hence $p_m = p_n$ (since both are prime), so $n = m$.
	In this branch the \emph{only} possible solution at modulus $s$ is
	$n = m$, and this case is excluded for $m > 1$ by the following
	argument.
	
	If $n = m$ were a solution it would require $2m \mid 2p_m^2$,
	i.e., $m \mid p_m^2$.
	Since $p_m$ is prime, $m \mid p_m^2$ implies $m \in \{1, p_m, p_m^2\}$.
	The case $m = p_m^2$ is impossible for $m \geq 2$: since $p_m \geq 2$
	we have $p_m^2 \geq 2p_m > 2m > m$, a contradiction.
	The case $m = p_m$ is ruled out by an elementary observation: since
	$p_1 = 2, p_2 = 3, \ldots$ is a strictly increasing sequence of
	integers starting at~$2$, we have $p_m \geq m + 1 > m$ for all
	$m \geq 1$, so $p_m \neq m$.
	Hence $n = m$ is not a valid solution for any $m \geq 2$.
	For $m = 1$, by contrast, $n = m = 1$ is a valid solution: $m + n = 2$
	divides $p_1^2 + p_1^2 = 8$, so $a(1) = 1$.
	The algorithm handles this edge case before the main search.
	For $m \geq 2$, the $n = m$ branch is excluded and the algorithm
	proceeds to the main search loop.
	
	A classical result (see \cite{IrelandRosen}, Ch.~5) gives:
	
	\begin{proposition}\label{prop:qr}
		$-1$ is a quadratic residue modulo $s$ if and only if both of the
		following hold:
		\begin{enumerate}
			\item[(a)] $4 \nmid s$ \textup{(}equivalently, $v_2(s) \leq 1$,
			where $v_2(s)$ denotes the 2-adic valuation of $s$\textup{)}, and
			\item[(b)] no prime $q \equiv 3 \pmod{4}$ divides $s$.
		\end{enumerate}
	\end{proposition}
	
	\begin{remark}\label{rem:jacobi}
		Condition~(b) is strictly stronger than requiring the Jacobi symbol
		$\left(\frac{-1}{s}\right) = 1$.
		For odd $s$, the Jacobi symbol $\left(\frac{-1}{s}\right) = 1$ if and
		only if the \emph{sum} $E = \sum_{q \equiv 3 \pmod{4}} v_q(s)$ of the
		exponents of all primes $q \equiv 3 \pmod{4}$ dividing $s$ is even
		(including $E = 0$).
		By contrast, actual solvability of $x^2 \equiv -1 \pmod{s}$ requires
		that no such prime divides $s$ at all, i.e., $E = 0$.
		Thus the Jacobi symbol can equal~$1$ even when $-1$ is \emph{not} a
		quadratic residue modulo~$s$; a simple example is $s = 21 = 3 \times 7$,
		for which $E = 1 + 1 = 2$ is even so $\left(\frac{-1}{21}\right) = 1$,
		yet $-1$ is not a quadratic residue modulo $21$ because condition~(b)
		fails.
		The Jacobi symbol is therefore only a necessary condition for
		solvability, not a sufficient one.
		For $s \equiv 2 \pmod{4}$, the Jacobi symbol is undefined (it is only
		defined for odd positive denominators); solvability of
		$x^2 \equiv -1 \pmod{s}$ in this case requires that no odd prime
		$q \equiv 3 \pmod{4}$ divides $s$, which is exactly condition~(b)
		applied to the odd part of $s$.
	\end{remark}
	
	When $-1$ is a QR mod $s$ and $s$ is odd (so every prime divisor of
	$s$ is $\equiv 1 \pmod{4}$), there are exactly $2^{\omega(s)}$ square
	roots of $-1$ modulo $s$, where $\omega(s)$ denotes the number of
	distinct prime factors of~$s$ (see \cite{IrelandRosen}, Ch.~5,
	Theorem~9).
	When $s \equiv 2 \pmod{4}$ (i.e., $v_2(s) = 1$), condition~(a) is
	satisfied and the factor of~$2$ contributes exactly one solution to
	$x^2 \equiv -1 \pmod{2}$, so the total count is $2^{\omega(s)-1}$;
	this case is handled correctly by the filter
	(see Section~\ref{sec:filter}).
	
	The modular filter in our implementation checks condition~(a) exactly
	and condition~(b) partially (up to a trial-division bound); see
	Section~\ref{sec:filter} for details.
	Candidates not eliminated by the filter are tested by direct arithmetic.
	
	% ----------------------------------------------------------
	\section{Computational Methods}
	% ----------------------------------------------------------
	
	\subsection{Algorithm and modular filter}
	\label{sec:filter}
	
	Our implementation (written in C) uses a segmented sieve of
	Eratosthenes~\cite{CrandallPomerance} to build a table of all primes up
	to a specified limit, together with a modular filter to skip values of
	$n$ for which $s = m + n$ admits no square root of $-1$, before testing
	the divisibility condition directly.
	
	The filter checks, for each candidate $n$, two conditions:
	\begin{enumerate}
		\item[(a)] $4 \nmid s$, i.e., $s \not\equiv 0 \pmod{4}$
		(condition~(a) of Proposition~\ref{prop:qr}, checked exactly); and
		\item[(b$'$)] no prime $q \equiv 3 \pmod{4}$ up to a trial-division
		bound $B \approx 10^4$ divides $s$ (a partial check of
		condition~(b) of Proposition~\ref{prop:qr}).
	\end{enumerate}
	This is a safe filter: it eliminates only values of $s$ for which
	$x^2 \equiv -1 \pmod{s}$ provably has no solution, so no valid
	candidate is ever discarded.
	More precisely, if either condition~(a) or~(b) of
	Proposition~\ref{prop:qr} is violated for~$s$, then no
	$x \in \mathbb{Z}$ satisfies $x^2 \equiv -1 \pmod{s}$, and
	consequently the divisibility $(m+n) \mid p_m^2 + p_n^2$ cannot hold
	when $\gcd(p_m, s) = 1$ (since it would require
	$(p_n p_m^{-1})^2 \equiv -1 \pmod{s}$); the filter is therefore
	\emph{sound} in the sense that it discards only provably impossible
	candidates.
	Values of $s$ that have a prime factor $q \equiv 3 \pmod{4}$ larger
	than $B$ are not eliminated by the filter; they pass through to the
	direct divisibility test, which may still fail, resulting in a small
	amount of redundant arithmetic.
	Across the range $m = 1, \ldots, 100\,000$, the filter eliminates
	empirically $73$--$77\%$ of candidate values of $n$ before any
	arithmetic is performed.
	Numerically, approximately $83\%$ of integers near $10^9$ possess at
	least one prime factor $q \equiv 3 \pmod{4}$ to an \emph{odd} power
	(a consequence of the Landau--Ramanujan asymptotic
	formula~\cite{Landau1908} evaluated at finite scale).
	The filter's condition is actually \emph{broader}: it checks for
	divisibility by any prime $q \equiv 3 \pmod{4}$ to \emph{any} power
	(odd or even), which defines a superset of the ``odd power'' set;
	the theoretical elimination rate for the filter's exact condition is
	therefore slightly \emph{higher} than $83\%$.
	The observed rate of $73$--$77\%$ is lower than this theoretical
	expectation solely because the trial-division bound $B \approx 10^4$
	leaves large prime factors unchecked.
	
	All integer arithmetic uses 64-bit signed integer types
	(\texttt{int64\_t}) throughout.
	The divisibility check $(p_m^2 + p_n^2) \bmod s = 0$ is implemented
	using explicit reduction before multiplication:
	\[
	\bigl((p_m^2 \bmod s) + (p_n \bmod s)^2\bigr) \bmod s \;=\; 0.
	\]
	The formula treats $p_m$ and $p_n$ asymmetrically for a concrete
	reason: for $m \leq 100\,000$ we have $p_m \leq 1\,299\,709$
	(the exact value of $p_{100\,000}$),
	so $p_m^2 \leq 1.69 \times 10^{12}$ fits directly in \texttt{int64\_t}
	and is reduced after squaring; by contrast, $p_n$ can reach
	$4.71 \times 10^{10}$, so $p_n^2 \approx 2.2 \times 10^{21}$ would
	overflow \texttt{int64\_t} unless $p_n$ is reduced \emph{before}
	squaring.
	After reduction, $\mathtt{pnr} = p_n \bmod s < s \leq m + n \leq
	2.0001 \times 10^9$, and hence
	\[
	\mathtt{pnr}^2 < (2.0001 \times 10^9)^2
	= 4.0004 \times 10^{18}
	< 2^{63} - 1 \approx 9.22 \times 10^{18}.
	\]
	To close the overflow argument: the sum
	$(p_m^2 \bmod s) + \mathtt{pnr}^2$ satisfies
	$(p_m^2 \bmod s) \leq s - 1 < 2.0001 \times 10^9$
	and $\mathtt{pnr}^2 \leq 4.0004 \times 10^{18}$,
	so the sum is at most $4.0004 \times 10^{18} + 2.0001 \times 10^9
	< 2^{63} - 1$, and the final reduction is safe.
	No overflow occurs on \texttt{int64\_t}; neither \texttt{\_\_int128}
	nor arbitrary-precision arithmetic is required.
	Prime generation uses a bitset segmented sieve of Eratosthenes
	storing one bit per odd number, which halves memory relative to a
	full-integer representation and is fully deterministic: every number
	up to the sieve limit is either confirmed prime or confirmed composite;
	no probabilistic primality test is employed.
	
	\subsection{Computation phases}
	
	The computation proceeded in phases:
	\begin{enumerate}
		\item \emph{Phases~1--2} ($m \leq 10\,000$, bound $n \leq 10^8$):
		established baseline results, resolving $9998$ of $10\,000$ cases.
		\item \emph{Phase~3} ($m \leq 10\,000$,
		bound $n \leq 7.6 \times 10^8$):
		resolved the remaining $2$ cases, confirming
		$a(4703) = 760\,027\,770$ (previously noted by Sun~\cite{Sun2019})
		and establishing $a(5263) = 420\,423\,462$.
		\item \emph{Phases~4--6} ($m \leq 100\,000$,
		bound $n \leq 2 \times 10^9$):
		extended the dataset to $m = 100\,000$, resolving $99\,951$ of
		$100\,000$ cases.
	\end{enumerate}
	
	All computations were performed on a personal workstation equipped with
	an AMD Ryzen~9 7940HS processor and 32\,GB of RAM\@.
	The final phase (Phase~6, $m \leq 100\,000$) required $597$ seconds of
	wall-clock time ($\approx 10$~minutes) and approximately $20.8$\,GB of
	RAM: $\sim 3.2$\,GB for the packed-bit sieve up to
	$5.14 \times 10^{10}$ (one bit per odd number) and $\sim 17.6$\,GB for
	the prime index table storing all $2\,176\,824\,767$ primes below
	$5.14 \times 10^{10}$ as 8-byte integers.
	The search bound $n \leq 2 \times 10^9$ corresponds to
	$p_n \leq 47\,055\,833\,459 \approx 4.71 \times 10^{10}$.
	The sieve is built to $5.14 \times 10^{10}$, a margin of
	approximately $9\%$ above $p_{2 \times 10^9}$, to provide a
	comfortable buffer that avoids edge effects and ensures all needed
	primes are available.
	
	\subsection{Independent verification}
	\label{sec:verification}
	
	Results were verified by three methods, as described below.
	Minimality of $a(m)$ is fully independently verified for $290$
	values; for all other values it depends on the correctness of the
	exhaustive C~search.
	
	\begin{enumerate}
		\item \emph{Cross-consistency}: all $10\,000$ values for
		$m \leq 10\,000$ agree exactly between Phase~3 and Phase~6 outputs,
		confirming implementation stability and determinism across
		independent runs.
		Note that this method confirms reproducibility of the algorithm,
		not its correctness; two runs of the same implementation will
		agree regardless of any systematic bug.
		\item \emph{Direct divisibility via \texttt{SymPy}}: for each of the
		$398$ exceptional resolved cases ($a(m) > 10^8$), the values $p_m$
		and $p_{a(m)}$ were computed independently using
		\texttt{sympy.prime()}~\cite{SymPy}, which implements its own
		prime-generation algorithm entirely separate from the C~sieve, and
		we verified $(p_m^2 + p_n^2) \bmod (m+n) = 0$ directly.
		All $398$ verified without error.
		The \texttt{SymPy} and Python scripts take as input only the values
		of $m$ and the claimed $a(m)$; they do not use the prime table
		produced by the C~program.
		This method confirms the divisibility condition for the claimed
		values but \emph{not} minimality; consequently, for the $398$
		exceptional resolved cases, minimality is not independently
		verified beyond the correctness of the C~implementation.
		\item \emph{Independent sieve}: an independent C implementation
		(\texttt{verify\_min.c}) using a segmented sieve of Eratosthenes,
		with no code shared with the Phase-6 program, verified $290$ cases
		by exhaustive search from $n = 1$, confirming both the divisibility
		condition \emph{and} the minimality of $a(m)$.
		The $290$ cases were selected in three groups:
		\begin{itemize}
			\item \emph{Group~B1}: the $100$ smallest exceptional cases
			($a(m) \in [1.00 \times 10^8,\; 1.52 \times 10^8]$), providing
			the first independent minimality check for exceptional values.
			\item \emph{Group~B2}: $150$ randomly selected non-exceptional
			cases ($a(m) < 10^8$, seed~$= 42$).
			\item \emph{Group~C}: all $40$ resolved cases with
			$a(m) > 10^9$, including $m = 14\,740$
			($a(m) = 1{,}994{,}463{,}433$) and $m = 16\,167$
			($a(m) = 1{,}881{,}095{,}843$), the two largest values in the
			dataset.
		\end{itemize}
		All $290$ cases were confirmed without discrepancy.
		The verification required $22.2$ minutes of wall-clock time on the
		author's workstation (AMD Ryzen~9 7940HS), with individual case
		times ranging from under $0.1$~seconds (small $a(m)$) to
		$40.2$~seconds ($m = 14\,740$, $a(m) \approx 2 \times 10^9$),
		consistent with the linear dependence of cost on $a(m)$.
	\end{enumerate}
	
	In total, $10\,688$ individual verifications were performed across all
	three methods (counting cases checked by multiple methods separately),
	with zero discrepancies.
	
	% ----------------------------------------------------------
	\section{Results}
	% ----------------------------------------------------------
	
	\subsection{Summary}
	
	For $m = 1, \ldots, 100\,000$, we determined $a(m)$ in $99\,951$ cases.
	The remaining $49$ values are listed and discussed in
	Section~\ref{sec:notfound}.
	Table~\ref{tab:top10} lists the ten largest values of $a(m)$ found.
	
	We call a resolved $m$ \emph{exceptional} if $a(m) > 10^8$.
	This threshold is approximately $1500$ times the median value of $a(m)$
	($\approx 65\,000$, computed from the $99\,951$ resolved values).
	There are $E_{\mathrm{res}}(10^5) = 398$ such resolved exceptional
	cases.
	Since the $49$ unresolved values all satisfy $a(m) > 2 \times 10^9$
	by definition of the search bound, the true exceptional count
	satisfies $E(10^5) \geq 447$, where
	\[
	E(M) \;=\; E_{\mathrm{res}}(M)
	+ \#\bigl\{m \leq M : m \text{ unresolved}\bigr\}
	\]
	denotes the total exceptional count; the quantity $E_{\mathrm{res}}(M)$
	counts only resolved exceptional cases throughout this paper.
	
	The complete dataset (all $99\,951$ resolved values and the list of
	$49$ unresolved~$m$) will be submitted to the OEIS as an extension of
	the b-file for A247975.
	
	\begin{table}[ht]
		\centering
		\caption{Ten largest values of $a(m)$ for $m \leq 100\,000$.}
		\label{tab:top10}
		\begin{tabular}{rrrr}
			\toprule
			$m$ & $a(m)$ & $p_m$ & $a(m)/(m \log m)$ \\
			\midrule
			14\,740 & 1\,994\,463\,433 &  160\,669 & 14\,097 \\
			77\,673 & 1\,985\,513\,552 &  988\,541 &  2270 \\
			16\,167 & 1\,881\,095\,843 &  178\,093 & 12\,007 \\
			81\,709 & 1\,876\,857\,276 & 1\,044\,167 &  2031 \\
			99\,544 & 1\,796\,084\,178 & 1\,293\,169 &  1568 \\
			76\,412 & 1\,780\,612\,094 &  970\,861 &  2072 \\
			43\,197 & 1\,779\,577\,013 &  521\,753 &  3860 \\
			96\,127 & 1\,729\,717\,251 & 1\,245\,479 &  1568 \\
			26\,193 & 1\,659\,619\,897 &  302\,483 &  6228 \\
			66\,932 & 1\,612\,520\,046 &  840\,643 &  2168 \\
			\bottomrule
		\end{tabular}
	\end{table}
	
	\subsection{Outstanding cases}
	
	The five largest Cram\'er-type ratios $a(m)/(m\log m)$ among resolved
	$m \leq 100\,000$ are listed in Table~\ref{tab:ratio}.
	Since $p_m \sim m\log m$ by the prime number theorem
	(see \cite{Davenport}, Ch.~18), the ratio
	$C(m) = a(m)/(m\log m)$ is approximately $a(m)/p_m$, normalising the
	solution index by the size of the prime $p_m$.
	Under Cram\'er's probabilistic model for primes~\cite{Cramer1936},
	treating primes as independent random events with local density
	$1/\log n$ near $n$, the expected index at which one encounters a prime
	satisfying a given congruence condition of modulus of order $m$ is of
	order $m$; hence $C(m)$ measures how many times larger the actual
	solution is compared to this heuristic baseline.
	Notably, $m = 4703$ retains the highest ratio ($\approx 19\,111$)
	across the entire range $m \leq 100\,000$.
	
	\begin{table}[ht]
		\centering
		\caption{Top five resolved values by Cram\'er-type ratio
			$a(m)/(m\log m)$.}
		\label{tab:ratio}
		\begin{tabular}{rrrrr}
			\toprule
			Rank & $m$ & $a(m)$ & $p_m$ & $a(m)/(m\log m)$ \\
			\midrule
			1 &  4703 &  760\,027\,770 &  45\,329 & 19\,111 \\
			2 & 14\,740 & 1\,994\,463\,433 & 160\,669 & 14\,097 \\
			3 & 16\,167 & 1\,881\,095\,843 & 178\,093 & 12\,007 \\
			4 &  5263 &  420\,423\,462 &  51\,437 &  9323 \\
			5 & 21\,907 & 1\,590\,716\,535 & 248\,323 &  7265 \\
			\bottomrule
		\end{tabular}
	\end{table}
	
	For $m = 4703$: $p_m = 45\,329$,
	$p_{a(m)} = 17\,111\,249\,191$ (independently verified via \texttt{SymPy}
	and WolframAlpha),
	$s = m + a(m) = 760\,032\,473$.
	For $m = 14\,740$: $p_m = 160\,669$,
	$p_{a(m)} = 46\,919\,734\,859$,
	$s = 1\,994\,478\,173$.
	
	\subsection{Growth of resolved exceptional cases}
	
	Let $E_{\mathrm{res}}(M)$ denote the number of resolved values
	$m \leq M$ with $a(m) > 10^8$.
	Table~\ref{tab:growth} shows the growth of $E_{\mathrm{res}}(M)$ over
	successive ranges of $10\,000$.
	
	\begin{table}[ht]
		\centering
		\caption{Count of resolved exceptional cases $a(m) > 10^8$
			by range of $m$.}
		\label{tab:growth}
		\begin{tabular}{lr}
			\toprule
			Range of $m$ & Count \\
			\midrule
			$1$--$10\,000$   &   2 \\
			$10\,001$--$20\,000$ &  27 \\
			$20\,001$--$30\,000$ &  21 \\
			$30\,001$--$40\,000$ &  25 \\
			$40\,001$--$50\,000$ &  41 \\
			$50\,001$--$60\,000$ &  48 \\
			$60\,001$--$70\,000$ &  44 \\
			$70\,001$--$80\,000$ &  60 \\
			$80\,001$--$90\,000$ &  62 \\
			$90\,001$--$100\,000$ &  68 \\
			\midrule
			Total & 398 \\
			\bottomrule
		\end{tabular}
	\end{table}
	
	% ----------------------------------------------------------
	\section{Statistical Analysis}
	% ----------------------------------------------------------
	
	\subsection{Heavy-tailed distribution}
	
	The distribution of $a(m)$ for resolved $m = 1, \ldots, 100\,000$ is
	strongly heavy-tailed.
	We apply the Hill estimator~\cite{Hill1975} as an exploratory
	diagnostic tool; since the data are deterministic and not i.i.d.,
	this should be understood as descriptive rather than as a formal
	inference about an underlying distribution (see \cite{Clauset2009}
	for a discussion of the limitations of power-law fitting in empirical
	data).
	Given $N_s$ order statistics $X_{(1)} \leq \cdots \leq X_{(N_s)}$, the
	Hill estimator with threshold parameter $k$ is defined as
	\[
	\hat{\alpha}_{\mathrm{Hill}}(k)
	= \left(
	\frac{1}{k} \sum_{i=1}^{k}
	\log \frac{X_{(N_s-i+1)}}{X_{(N_s-k)}}
	\right)^{-1}.
	\]
	The Hill statistic, computed over the top $10\%$ of resolved values
	($k = 9995$, threshold $a(m) > 1\,142\,143$), takes the value
	\[
	\hat{\alpha}_{\mathrm{Hill}} = 0.713.
	\]
	The statistic varies with threshold:
	\[
	\hat{\alpha}(1\%) = 0.817, \quad
	\hat{\alpha}(5\%) = 0.732, \quad
	\hat{\alpha}(10\%) = 0.713, \quad
	\hat{\alpha}(20\%) = 0.676.
	\]
	The threshold-dependence indicates that the tail is not a pure power
	law; the values of $\hat{\alpha}_{\mathrm{Hill}}$ should be interpreted
	only as an indication of heavy-tailed behaviour, not as evidence of a
	specific parametric form.
	
	\subsection{Cram\'er-type ratio}
	
	Define the Cram\'er-type ratio $C(m) = a(m)/(m \log m)$ as in
	Section~4.2.
	The median of $C(m)$ over resolved $m = 2, \ldots, 100\,000$
	(computed from the $99\,951$ resolved values) is
	\[
	\widetilde{C} \approx 0.157,
	\]
	stable and well-defined across the full range.
	
	% ----------------------------------------------------------
	\section{Heuristic Model and Comparison with Data}
	\label{sec:heuristic}
	% ----------------------------------------------------------
	
	\subsection{Density argument}
	
	Fix $m$ and consider the search for a valid $n$.
	For a given candidate $s = m + n$, a solution exists if and only if
	$p_n$ falls in one of the $2^{\omega(s)}$ residue classes of the form
	$r \cdot p_m \pmod{s}$, where $r$ ranges over the square roots of
	$-1 \pmod{s}$.
	By Dirichlet's theorem~\cite{Davenport}, the asymptotic density
	of primes in each such class is $\varphi(s)^{-1}$ for fixed $s$,
	where $\varphi(s)$ denotes Euler's totient function (the number of
	integers in $\{1, \ldots, s\}$ coprime to $s$); we note
	that applying this for moduli $s$ that grow up to $\sim 2 \times 10^9$
	while checking primes up to $\sim 4.7 \times 10^{10}$ goes beyond the
	range of standard uniformity results (Siegel--Walfisz,
	see \cite{Davenport}, Ch.~21), and the following heuristic should be
	treated as qualitative.
	The total proportion of admissible residue classes modulo $s$ is
	\[
	D(s) \;=\; \frac{2^{\omega(s)}}{\varphi(s)}.
	\]
	Note that $D(s)$ is the \emph{proportion of admissible residue classes},
	not a probability of success for a randomly chosen prime near $s$;
	the following derivation should be regarded as a rough heuristic
	envelope rather than a rigorous estimate.
	
	Using the normal-order estimate $\omega(s) \approx \log\log s$
	\cite{HardyRamanujan} and the extremal lower bound
	$\varphi(s) \gtrsim s \cdot e^{-\gamma}/\log\log s$
	\cite{Mertens}
	(the first is a normal-order result valid for almost all $s$, the
	second an extremal bound realised at primorials; their combination
	has no rigorous standing), one obtains:
	\[
	D(s) \;\lesssim\;
	\frac{(\log s)^{\log 2} \cdot e^{\gamma} \log\log s}{s}.
	\]
	Treating each candidate $n$ as independently selected with probability
	$D(m+n)$ of yielding a valid solution, the expected value of $a(m)$
	satisfies $\sum_{n=1}^{N} D(m+n) \approx 1$, which for slowly varying
	$D$ reduces to $\int_m^{m+N} D(s)\,ds \approx 1$.
	Substituting the bound on $D(s)$ and integrating yields
	\[
	\int_m^{m+N} \frac{(\log s)^{\log 2} \cdot e^{\gamma} \log\log s}{s}
	\,ds
	\;\approx\;
	e^{\gamma} \cdot (\log m)^{\log 2} \cdot \log\log m
	\cdot \frac{N}{m}
	\;\approx\; 1,
	\]
	where the approximation treats $\log s \approx \log m$ and
	$\log\log s \approx \log\log m$ over the range $[m, m+N]$, valid
	since $N \ll m$ in the regime of interest.
	This gives $N \sim m/(e^{\gamma} (\log m)^{\log 2} \log\log m)$;
	for $m = 10^5$ this yields
	$\approx 0.04\,m$ --- well below the observed median
	$a(m) \approx 1.8\,m$.
	This large discrepancy shows that the density model does not explain
	the observed scale of $a(m)$; the Cram\'er-type normalization
	$m\log m$ is motivated independently by the prime-index analogy
	(Section~5.2), not derived from $D(s)$.
	
	\subsection{Comparison with data}
	
	The density model provides a qualitative explanation for the heavy tail
	(the factor $(\log s)^{\log 2}$ in $D(s)$ produces slower-than-power-law
	decay), but fails quantitatively: it predicts the typical scale of
	$a(m)$ as $\sim m/(e^{\gamma} (\log m)^{\log 2} \log\log m)
	\approx 0.04\,m$, whereas the typical $a(m)$ scales as
	approximately $1.8\,m$.
	The growing count $E_{\mathrm{res}}(M)$ in Table~\ref{tab:growth} is
	qualitatively consistent with increasing exceptional-case frequency,
	but the growth factor of $199$ from $E_{\mathrm{res}}(10^4) = 2$ to
	$E_{\mathrm{res}}(10^5) = 398$ is far larger than any simple log-power
	formula would predict.
	We therefore regard the density model as providing only a qualitative
	motivation for heavy-tailed behaviour, and do not propose a
	quantitative formula for $E_{\mathrm{res}}(M)$.
	
	The outstanding Cram\'er-type ratio of $m = 4703$ (ratio $\approx 19\,111$,
	the global maximum over $m \leq 100\,000$) points toward richer arithmetic
	structure: specific properties of $p_{4703} = 45\,329$ and its
	interaction with the moduli $s = m+n$ may produce a systematic
	suppression of $D(s)$ not captured by the average model.
	A refined heuristic remains an open problem.
	
	% ----------------------------------------------------------
	\section{Unresolved Cases}
	\label{sec:notfound}
	% ----------------------------------------------------------
	
	The $49$ values of $m \leq 100\,000$ for which no solution was found
	within the search bound $n \leq 2 \times 10^9$ are listed in
	Table~\ref{tab:notfound}, together with $p_m$ and the filter-detected
	obstruction rate (percentage of candidate $n$ eliminated by the
	modular filter).
	
\begin{table}[ht]
	\centering
	\caption{The $49$ unresolved values of $m \leq 100\,000$, with $p_m$
		and filter-detected obstruction rate (\%).
		No solution was found with $n \leq 2 \times 10^9$.}
	\label{tab:notfound}
	\small
	\begin{tabular}{rrr@{\qquad}rrr}
		\toprule
		$m$ & $p_m$ & Obstr.\ (\%) & $m$ & $p_m$ & Obstr.\ (\%) \\
		\midrule
		11\,924 &  127\,289 & 75.1 & 68\,744 &  865\,327 & 76.9 \\
		18\,241 &  203\,227 & 75.5 & 69\,546 &  876\,191 & 76.9 \\
		19\,234 &  215\,317 & 75.6 & 71\,049 &  896\,723 & 76.9 \\
		20\,686 &  233\,173 & 75.7 & 72\,568 &  917\,759 & 76.9 \\
		22\,186 &  251\,809 & 75.7 & 76\,761 &  975\,497 & 77.0 \\
		22\,802 &  259\,577 & 75.7 & 76\,868 &  977\,087 & 77.0 \\
		23\,995 &  274\,453 & 75.8 & 81\,986 & 1\,047\,971 & 77.0 \\
		33\,566 &  396\,239 & 76.2 & 82\,096 & 1\,049\,519 & 77.0 \\
		36\,303 &  431\,833 & 76.2 & 82\,300 & 1\,052\,459 & 77.0 \\
		36\,726 &  437\,083 & 76.2 & 84\,521 & 1\,083\,371 & 77.1 \\
		37\,249 &  443\,837 & 76.2 & 86\,557 & 1\,111\,393 & 77.1 \\
		37\,253 &  443\,873 & 76.2 & 90\,005 & 1\,159\,597 & 77.1 \\
		37\,949 &  453\,181 & 76.3 & 90\,161 & 1\,161\,757 & 77.1 \\
		43\,297 &  523\,007 & 76.4 & 92\,335 & 1\,191\,563 & 77.1 \\
		48\,834 &  596\,419 & 76.5 & 93\,790 & 1\,212\,347 & 77.2 \\
		49\,889 &  610\,469 & 76.5 & 93\,890 & 1\,213\,651 & 77.1 \\
		51\,816 &  636\,277 & 76.5 & 94\,429 & 1\,221\,131 & 77.2 \\
		53\,963 &  665\,179 & 76.7 & 94\,768 & 1\,225\,981 & 77.2 \\
		55\,468 &  685\,453 & 76.7 & 97\,251 & 1\,260\,901 & 77.2 \\
		55\,874 &  690\,953 & 76.7 & 97\,351 & 1\,262\,453 & 77.2 \\
		61\,007 &  760\,367 & 76.8 & 98\,008 & 1\,271\,383 & 77.2 \\
		66\,257 &  831\,301 & 76.9 & 98\,258 & 1\,275\,011 & 77.3 \\
		67\,203 &  844\,253 & 76.9 & 98\,379 & 1\,276\,747 & 77.2 \\
		67\,273 &  845\,197 & 76.9 & 99\,105 & 1\,286\,953 & 77.2 \\
		67\,447 &  847\,621 & 76.9 &       &         &      \\
		\bottomrule
	\end{tabular}
\end{table}

All $49$ cases exhibit a filter-detected obstruction rate between
$75.1\%$ and $77.3\%$, somewhat higher than the typical $73$--$74\%$
for resolved cases (computed as the median obstruction rate over all
resolved $m \leq 100\,000$).
The difference of $1$--$2$ percentage points means that only about
$24\%$ of candidates survive the filter, compared to about $26\%$ for
typical resolved cases; the expected waiting time (inverse of the
survival rate) is therefore about $8\%$ larger, which is far too small
on its own to explain values exceeding $2 \times 10^9$.

To assess whether the $49$ unresolved cases are simply the extreme
tail of the $a(m)$ distribution, we compare them with the $49$
resolved cases having the largest $a(m)$.
The top~$49$ resolved cases have
$a(m) \in [846\,007\,757,\ 1\,994\,463\,433]$ with median
$\approx 1.23 \times 10^9$; only $2$ of these exceed
$1.9 \times 10^9$, and none reaches $2 \times 10^9$.
By contrast, all $49$ unresolved cases satisfy $a(m) > 2 \times 10^9$
by definition of the search bound.
The gap between the largest resolved value ($1\,994\,463\,433$ at $m = 14\,740$)
and the search bound suggests that the unresolved cases are not merely
a statistical continuation of the resolved tail, but may possess an
additional structural property --- consistent with their slightly
elevated obstruction rates.
Whether this reflects a genuine arithmetic distinction or an extreme
statistical fluctuation remains an open question.

We cannot rule out that some of these $49$ cases represent genuine
counterexamples to Conjecture~\ref{conj:sun}, though this seems
unlikely given that the resolved cases $m = 14\,740$ and $m = 16\,167$
required $a(m) \approx 1.9 \times 10^9$ and were themselves unresolved
in earlier phases.
The arithmetic characterization of the high-obstruction set remains an
open problem.

% ----------------------------------------------------------
\section{Conclusion}
% ----------------------------------------------------------

We have determined $a(m)$ for $99\,951$ values of $m \leq 100\,000$,
with $49$ values remaining unresolved within our search bound
$n \leq 2 \times 10^9$.
The resolved dataset reveals a heavy-tailed distribution with Hill
statistic $\hat{\alpha}_{\mathrm{Hill}} \approx 0.713$, a growing
density of exceptional cases, and a persistent outstanding case at
$m = 4703$ whose Cram\'er-type ratio remains the global maximum over
the entire range studied.

The heuristic density model provides a qualitative explanation for the
heavy tail but fails quantitatively to predict the scale of $a(m)$ or
the growth rate of exceptional cases.
Natural directions for future work include the arithmetic
characterization of the high-obstruction cases, an extended
computational search for the $49$ unresolved values, and a refined
probabilistic model for the distribution of $a(m)$
(see also \cite{Granville1995} for a discussion of corrections to
the classical Cram\'er heuristic).

\paragraph{Data availability.}
The complete dataset ($99\,951$ resolved values and $49$ unresolved
$m$-values) will be submitted to the OEIS as an extension of the
b-file for A247975.
The C~source code for the main search (\texttt{sun41\_v6b.c}), the
independent verification sieve (\texttt{verify\_min.c}), the
Python/SymPy verification scripts, and the full machine-readable
dataset are publicly available at:
\begin{itemize}
  \item GitHub repository: \url{https://github.com/carcorti/A247975}
  \item Permanent Zenodo archive (DOI):
        \url{https://doi.org/10.5281/zenodo.18920372}
\end{itemize}

% ----------------------------------------------------------
\section*{Acknowledgments}
% ----------------------------------------------------------

The author thanks Zhi-Wei Sun for the beautiful conjecture that
motivated this work.
The computational implementation, including the C~search program and
verification sieve, was developed with extensive assistance from
Anthropic's Claude language models (\texttt{claude-sonnet-4-6} and
\texttt{claude-opus-4-6}); Claude also assisted substantially in the
statistical analysis and in the preparation of the manuscript.
All computations were executed and independently verified by the author,
who bears full scientific responsibility for the results.
The author also thanks Chai Wah Wu, whose extension of the OEIS b-file
for A247975 to $m \leq 10\,000$ provided a valuable cross-check for
our results.

% ----------------------------------------------------------
\begin{thebibliography}{99}
	
	\bibitem{Sun2019}
	Zhi-Wei Sun,
	\textit{Some new problems in additive combinatorics},
	Nanjing Univ.\ J.\ Math.\ Biquarterly \textbf{36} (2019), no.~2,
	134--155; preprint: arXiv:1309.1679 [math.NT] (submitted 2013,
	v9: 2020).
	Conjecture~4.1(i) appears on p.~15; Remark~4.1 notes the case
	$m = 4703$. The conjecture was submitted to the OEIS in
	September~2014.
	\url{https://arxiv.org/abs/1309.1679}
	
	\bibitem{OEIS_A247975}
	OEIS Foundation,
	\textit{Sequence A247975},
	The On-Line Encyclopedia of Integer Sequences.
	\url{https://oeis.org/A247975}
	
	\bibitem{Hill1975}
	B.~M.~Hill,
	\textit{A simple general approach to inference about the tail of a
		distribution},
	Ann.\ Statist.\ \textbf{3} (1975), 1163--1174.
	
	\bibitem{IrelandRosen}
	K.~Ireland and M.~Rosen,
	\textit{A Classical Introduction to Modern Number Theory},
	2nd~ed., Springer, New York, 1990.
	Ch.~5, Theorem~9 for quadratic residues modulo composite integers.
	
	\bibitem{Davenport}
	H.~Davenport,
	\textit{Multiplicative Number Theory},
	3rd~ed., Springer, New York, 2000.
	Ch.~4 for Dirichlet's theorem; Ch.~21 for the Siegel--Walfisz
	theorem and uniformity in arithmetic progressions.
	
	\bibitem{SymPy}
	A.~Meurer et al.,
	\textit{SymPy: symbolic computing in Python},
	PeerJ Computer Science \textbf{3} (2017), e103.
	Version~1.12 used for independent prime computation.
	
	\bibitem{HardyRamanujan}
	G.~H.~Hardy and S.~Ramanujan,
	\textit{The normal number of prime factors of a number~$n$},
	Quart.\ J.\ Math.\ \textbf{48} (1917), 76--92.
	
	\bibitem{Mertens}
	F.~Mertens,
	\textit{Ein Beitrag zur analytischen Zahlentheorie},
	J.\ Reine Angew.\ Math.\ \textbf{78} (1874), 46--62.
	See also G.~Tenenbaum,
	\textit{Introduction to Analytic and Probabilistic Number Theory},
	Cambridge Univ.\ Press, 1995, Ch.~I.3.
	
	\bibitem{Cramer1936}
	H.~Cram\'er,
	\textit{On the order of magnitude of the difference between consecutive
		prime numbers},
	Acta Arith.\ \textbf{2} (1936), 23--46.
	
	\bibitem{WuOEIS}
	Chai Wah Wu,
	\textit{Table of $n$, $a(n)$ for $n = 1..10\,000$},
	contribution to OEIS sequence A247975, 2014.
	\url{https://oeis.org/A247975/b247975.txt}
	
	\bibitem{Clauset2009}
	A.~Clauset, C.~R.~Shalizi, and M.~E.~J.~Newman,
	\textit{Power-law distributions in empirical data},
	SIAM Review \textbf{51} (2009), 661--703.
	
	\bibitem{CrandallPomerance}
	R.~Crandall and C.~Pomerance,
	\textit{Prime Numbers: A Computational Perspective},
	2nd~ed., Springer, New York, 2005.
	Ch.~3 for sieve algorithms.
	
	\bibitem{Landau1908}
	E.~Landau,
	\textit{\"Uber die Einteilung der positiven ganzen Zahlen in vier
		Klassen nach der Mindestzahl der zu ihrer additiven Zusammensetzung
		erforderlichen Quadrate},
	Arch.\ Math.\ Phys.\ \textbf{13} (1908), 305--312.
	Establishes the asymptotic count ${\sim}\,Kx/\sqrt{\ln x}$
	(with $K \approx 0.7642$) of integers up to $x$ with no prime factor
	$q \equiv 3 \pmod{4}$ to odd power; the natural density of this
	set is zero.
	The value ${\sim}83\%$ cited in Section~3.1 is a numerical consequence
	of this formula evaluated at the finite scale $x \approx 10^9$,
	not a direct theorem prediction.
	
	\bibitem{Granville1995}
	A.~Granville,
	\textit{Harald Cram\'er and the distribution of prime numbers},
	Scand.\ Actuarial J.\ (1995), no.~1, 12--28.
	Provides a refined probabilistic model correcting the classical
	Cram\'er heuristic.
	
\end{thebibliography}

\end{document}
